% Source: https://artofproblemsolving.com/community/c3h2323115_arr0ws_firing_range
% Stable identity: aops-topic-2323115-q11
Let $y=\sqrt{5-x}\ge0$. Then the equation says $x=\sqrt{5-y}$, so
\[x^2=5-y,\qquad y^2=5-x,\qquad x>0,\ y\ge0.\]
Subtracting, $x^2-y^2=x-y$, so $(x-y)(x+y-1)=0$.
\begin{itemize}
\item If $x+y=1$, then $x^2=5-(1-x)=4+x$, so $x=\dfrac{1\pm\sqrt{17}}2$; the root with $x>0$ is about $2.56$, but then $y=1-x<0$, impossible.
\item If $x=y$, then $x^2+x-5=0$, so $x=\dfrac{-1+\sqrt{21}}{2}$ (the positive root).
\end{itemize}
Check: $x\approx1.791$, $5-x\approx3.209$, $\sqrt{3.209}\approx1.791$, and $\sqrt{5-1.791}\approx1.791$. So $\boxed{x=\dfrac{\sqrt{21}-1}{2}}$.
